\chapter*{Preface}
\addcontentsline{toc}{chapter}{Preface}

This book grows with the GPU port of the Weather Research and Forecasting model (WRF) and its fire models. It has
three aims:
\begin{enumerate}
  \item to explain WRF and WRF-Fire as they are implemented: the equations, the numerical methods, and how they are
    coded, routine by routine;
  \item to explain how such a code can be moved to graphics processors without changing a single bit of its
    results, and what that costs;
  \item to explain how the port is then made fast on NVIDIA A100 and H100 GPUs: what limits each kernel, and which
    changes help without changing the results.
\end{enumerate}

The book is written alongside the port. The chapters on the model describe WRF~4.6.0 as it is. The chapters on the
port describe the method and the code as far as the port has come. Measured results enter as the port produces
them; until then they are marked \tbr.

\paragraph{Conventions.}
\begin{itemize}
  \item A routine is cited with its file in the WRF source tree, for example
    \src{dyn_em/module_small_step_em.F}{advance_w}. Unless stated otherwise, the code is WRF~4.6.0 as imported into
    the repository, with the changes of the reproducible CPU reference build described in \cref{ch:repro}.
  \item Kernel names such as K-AW and test names such as T-TRACE refer to the port's plan (\code{plan.md}). The
    kernel catalogue is in \cref{app:kernels}, the test catalogue in \cref{app:tests}.
  \item Every equation names the routine that implements it, so that the reader can go from the mathematics to the
    code and back. Symbols are listed in \cref{app:notation}.
  \item Algorithms are written in pseudocode that keeps the loop order and the order of operations of the Fortran
    source, because on a GPU that order is what keeps results identical.
\end{itemize}

\paragraph{Sources.} The model description follows the WRF Technical Note \citep{skamarock2019} and the papers
cited in each chapter, checked against the code. Where the code and a paper differ, the book describes the code and
says so. The bibliography's entries are marked in \code{refs.bib} as checked against the publisher or not yet
checked.
